% !TeX root = ../../tesi.tex

\section{Risultati}

\subsection{Il prodotto dal punto di vista dell’utente}

Presento ChatCart dal punto di vista di chi lo utilizza, attraverso una sequenza di schermate: un
cliente che cerca un prodotto, lo confronta con altri, lo aggiunge al carrello e raggiunge il
pagamento; poi il commerciante che consulta il \emph{funnel} e le conversazioni nella
\emph{dashboard}. Dichiaro inoltre i limiti noti del prototipo.

\subsection{Risultati quantitativi}
\label{sec:quantitativi}

Riporto in tabelle i dati quantitativi del progetto: la copertura dei requisiti definiti nella
\secref{sec:requisiti}, la copertura dei \emph{test} con le scelte che la motivano, le righe di
codice per ciascuna parte dell’applicazione e i documenti che ho consegnato al \emph{team}.
